\section{La clase en el hijo de tamaño \texorpdfstring{\(j\)}{j}}
Fijemos la banda de [13, §5.4], la entrada regular en \(M=\lceil\delta n\rceil\), la difusión
\(C\), la brecha \(\gamma\) y las constantes del coste exterior y de salida \(c_*>2\beta\).
En \(H_j\) la historia ha permanecido regular hasta el hijo de tamaño \(j=\lfloor bn\rfloor\)
inclusive. Cada línea diagonal tiene menos de \(\lceil\alpha j\rceil\) casillas disponibles.
Denotemos por \(A\) su máscara binaria. Se usan tres matrices en este mismo soporte.
El punto de reinas \(p\) cumple \(j\max p_e\le C\) y tiene brecha de incidencia de filas y columnas
al menos \(\gamma\). La masa \(p\) de cada diagonal es a lo sumo \(C\alpha\) por integralidad.

El punto estacionario de torres \(P\) satisface en el soporte
\(|P_e/p_e-1|\le K_T(C\alpha+1/j)\) por la contracción en el tangente de
\pruebaen{D04}. Con \(C\alpha<\epsilon_T\), \(K_TC\alpha\le1/4\) y \(K_T/j\le1/4\)
se obtiene \(p/2\le P\le 3p/2\). En particular \(j\max P_e\le 2C\) y
\(\epsilon_A:=\max_f\sum_{e\in f}P_e\le2C\alpha\) en ambas familias diagonales.
La forma de incidencia de \(P\) domina la mitad de la de \(p\) en el complemento
del modo nulo \((\mathbf1,-\mathbf1)\). Sus autovalores son \(1\pm\sigma_i(P)\).
Así \(\|P-J_j\|_{\rm op}\le1-\gamma_T\), con
\(J_j=\mathbf1\mathbf1^t/j\) y \(\gamma_T=\min(1,\gamma)/2\).

Sea \(D\) el escalado bistocástico exacto de \(A\). La estacionariedad da
\(P_e(1-P_e)=a_rb_c\) en el soporte. Balancear \(P\) con el campo
\(\log(1-P_e)=O(1/j)\) da, por el inverso del jacobiano de márgenes en el calibre de
\pruebaen{N02}, \(D_e/P_e=1+O(1/j)\). La unicidad del escalado lo identifica con \(D\).
El error tiene sumas absolutas por fila y por columna de orden \(1/j\). La desigualdad
\(\|D-P\|_2\le\sqrt{\|D-P\|_1\|D-P\|_\infty}\) da la misma cota en norma de operador.
Para \(j\) grande pueden reservarse \(j\max D_e\le4C\) y
\(\|D-J_j\|_2\le1-\gamma_T/2\). Hall, aplicado a \(D\), da \(\per A>0\).

\paragraph{Orden de elección.}
Los márgenes de clase fijan primero los radios de balanceo y la fracción máxima de
supresiones \(a_*\), y luego \(\varepsilon\) para el orden común
\(Q_*=\lfloor\varepsilon j\rfloor\). Para la norma exacta se reserva el factor exponencial
por sitio de las prescripciones largas. Se escogen después la fracción de corte de Taylor,
la de prescripciones cortas y el umbral \(\eta_F\) en \pruebaen{V03}.
Sea \(\eta=\min(\eta_F,\eta_\Lambda)\), con el umbral canónico de \pruebaen{N03}.
Independientemente, los mismos márgenes fijan los radios locales y la base combinatoria
\(C_1\) con alcance \(8\log n\). Fíjense \(0<a<1\) con
\(a\log(\max(2,C_1))<1/4\) y después \(\rho\ge2\) con \(a\log\rho>3\).
Se toma \(\alpha<\alpha_{\rm hist}\) para satisfacer las condiciones de torres anteriores y
\(2C\alpha<\eta/(2\rho^2)\), reduciéndola para los márgenes de carga del exterior en \(D\),
si se expresan en esa tabla. Para \(j\) grande, \(\rho^2(\epsilon_A+1/j)\le\eta\).

Se reduce \(\delta_*\) para ambas series geométricas de B.1, se elige el corte
\(0<b<\min(\delta/2,\delta_*)\) con su condición de prescripción y se hace \(n\) grande
a \(b\) fijo. La constante de prescripción procede de la prueba de [12, teorema 8.1]
en la clase histórica. En particular se fija antes de \(b\). Se obtiene
\(j\ge\lceil\sqrt n\rceil\), \(C'(4v+12)\le j^{1/3}\) y \(v^2/j\le\delta_*\)
para \(v\le8\log n\). Estas son condiciones sobre \(n_0(b)\).

\paragraph{Activación por sitios.}
En una evaluación \(Y\subseteq A\), se mantienen las restricciones de filas y columnas
y se activan los conflictos diagonales entre sitios de \(Y\). Sea \(Z(Y)\) la probabilidad
de evitarlos bajo el matching uniforme de \(A\), \(F(Y)=\log Z(Y)\) una vez demostrada
la positividad, y \(\Lambda(Y)\) el cambio del predictor respecto de torres.
Ponemos \(e=F-\Lambda\) y
\(\Delta_U f=\sum_{Y\subseteq U}(-1)^{|U|-|Y|}f(Y)\).
Los tres resultados se prueban sobre máscaras fijas. La probabilidad de torres evalúa
\(F\) dentro de cada máscara. La posterior integración de B.1 usa el indicador \(H_j\)
bajo la ley original de historias \(\nu\).
